\section{Controlled red-team platforms and judge validity}\label{sec:C28:platforms}
DTap supplies controllable simulated environments, adversarial exploration and policy-based evaluation for agent security~\cite{corpus260504808}. The important contribution is an execution and measurement substrate that makes attack surfaces and outcomes inspectable. Simulations can support broad controlled comparisons, but their external validity depends on whether tools, policies and workflows capture the deployment conditions of interest.

A policy-based judge is still an evaluator. It can misclassify an action, omit a side effect, or accept a nominally blocked operation whose downstream consequences violate the intent. Judge validation should include positive, negative and ambiguous cases with authoritative fixture state. If an attack targets the judge or its observation path, the evaluation needs a separate trusted record of effects. Otherwise a successful deception can be scored as a failed attack.

\subsection{Attacker--defender self-play changes the data distribution}\label{sec:C28:selfplay}
GPT-Red's 2026 report trains an attacker against a population of concurrently improving defenders. It samples an environment and permitted insertion point, allows bounded defender queries, validates the submitted attack, then resumes the defender and updates both roles from their respective rewards. Invalid attacks receive zero reward. The defender reward includes legitimate-task completion, so the training is not generally a zero-sum refusal game~\cite{wallace2026gptred}. The official July 15 announcement links the report; the arXiv version was submitted July 28. These are separate publication dates, not evidence of an additional model revision.

For a general analysis, let $D$ be a probability measure on a controlled environment space $\mathcal E$, $z\sim D$ an environment, $u\sim\alpha_\phi(\cdot\mid z)$ a proposed modification in a candidate space $\mathcal U$, and $\xi\sim P_{\theta}(\cdot\mid z,u)$ a defender trajectory in $\Xi$. Let $v(z,u)\in\{0,1\}$ denote validity under the declared attacker access, $L(z,u,\xi)\in[0,1]$ the prohibited-effect score, and $U(z,\xi)\in[0,1]$ legitimate utility. Distinct objectives are
\begin{align}
 J_A(\phi;\theta)&=\mathbb E[v(z,u)L(z,u,\xi)],\\
 J_D(\theta;\phi)&=\mathbb E[U(z,\xi)-\lambda L(z,u,\xi)],
 \qquad\lambda\geq0.
\end{align}
This formulation is an analytical comparison, not the proprietary training objective. It explains why attacker success alone cannot judge the defender: refusing every operation can suppress $L$ while destroying $U$. A claimed equilibrium guarantee additionally needs a specified policy class, payoff, valid-action domain and convergence argument. Alternating stochastic updates in a partially observed agent environment do not automatically meet those conditions.

Co-adaptation makes the training curve ambiguous. Falling attack success can mean a stronger defender, a weaker or collapsed attacker, an easier sampled task mix, or a changed validity checker. Preserve frozen attacker and defender checkpoints and evaluate a cross-play matrix $R_{ij}=\mathbb E[L\mid\alpha_{\phi_i},\pi_{\theta_j},D_0]$ on a fixed held-out distribution $D_0$. Match query budgets and include benign utility, invalid-attack frequency, evaluator error and total generation cost. GPT-Red reports population diversity and held-out assessment; that evidence does not certify robustness to every unobserved tool, harness, or attacker~\cite{wallace2026gptred}. Feedback validity is especially important when a learned judge supplies both training reward and reported success.

\subsection{A modular adversarial interface is a measurement contract}\label{sec:C28:superred}
The inspected SuperRed revision separates targets, optimizers, security claims, observable events and controllable surfaces. \path{Optimizer.on_event} receives routed events; \texttt{Controller.run} orchestrates interaction; hierarchical \texttt{SecurityDomainTag} scopes constrain access~\cite{superred2026}. These are source-level interfaces, not proof that every wrapper exposes a complete trajectory. Its documented target lifecycle preserves durable per-task memory across ephemeral resets, and target-model spend is outside the optimizer's token budget. Consequently, a matched optimizer budget does not establish matched total experiment cost.

A defensive integration should first run a no-modification optimizer to check that wrapping preserves benign behavior, then verify access isolation using synthetic surfaces with known scope. Reset tests must distinguish new attempts against one evolving memory from independent clean tasks. A wrapper that inadvertently carries poisoned memory into a supposedly independent trial changes both the threat model and the uncertainty unit. Pin the controller and target adapter together: a revision to which events are observable can strengthen an attacker without changing its checkpoint. No SuperRed attack or training run is executed here.

Attack surfaces should be varied one at a time before studying combinations. Document injection, tool-result poisoning, skill modification and memory contamination can interact, but a combined successful attack does not identify which boundary failed. Use controlled ablations and counterfactual replay to test the causal chain. Preserve the benign task so the defender must continue useful work rather than simply classify a suspicious input in isolation.

Red-team artifacts should remain bounded and reproducible. Use isolated fixtures, synthetic confidential records and controlled destinations. The scientific record needs attacker access, permitted effects, payload category, budget, defender configuration and raw outcome evidence. It does not require exposing operational details that would enable misuse against live systems. A defensible survey can analyze the mechanism and controls without reproducing harmful deployment instructions.

